% !TeX root = ./rpz.tex
\chapter{{\fontsize{14pt}{21pt}\selectfont\MakeUppercase{Внедрение и эксплуатация}}}

\section{Анализ внедрения продукта}

В первой главе сформулированы цель разработки и перечень функциональных и
нефункциональных требований к системе.
В данном разделе проводится сопоставление этих требований с результатами реализации.

Все сформулированные функциональные требования реализованы в полном объeме.
Управление задачами (создание, редактирование, удаление с поддержкой мягкого удаления
через поле \texttt{is\_deleted}) реализовано на серверной стороне через REST API
и представлено в пользовательском интерфейсе в виде карточек с соответствующими
элементами управления.
Хранение географических координат обеспечивается полями \texttt{latitude} и
\texttt{longitude} в таблице \texttt{tasks} базы данных PostgreSQL.
Геокодирование адресов выполняется через Яндекс Карты JavaScript API~2.1:
реализован вывод до пяти предложенных вариантов с интерактивным выбором пользователя.
Матрица расстояний формируется с привлечением сервиса OSRM~\cite{osrm}; полученные значения
сохраняются в таблице \texttt{api\_distance\_cache} для повторного использования.
Алгоритм оптимизации, реализованный на серверной стороне, вычисляет оптимальный
порядок посещения точек с учeтом временных окон, длительности выполнения задач,
фиксированных начальной и конечной точек, а также попарных ограничений порядка,
рассчитывая для каждой остановки время прибытия, начала и окончания обслуживания.
Результат визуализируется на карте в виде полилинии с пронумерованными метками
(пресет \texttt{islands\#blueCircleIcon} Яндекс Карт), а также отображается в
упорядоченном списке задач.
Дополнительно реализована возможность ручного упорядочивания задач методом
перетаскивания посредством библиотеки \texttt{react-dnd}.
Экспорт задач реализован в двух форматах: CSV (библиотека \texttt{papaparse})
и Excel~XLSX (библиотека \texttt{xlsx}), а также доступен возврат
полной структуры маршрута в формате JSON через эндпоинт \texttt{GET /api/routes/:id}.
Дополнительно реализован импорт задач из файлов CSV и XLSX
с построчной валидацией через диалог импорта.

Нефункциональные требования также выполнены.
Серверная часть реализована на языке Go с использованием HTTP-фреймворка Gin.
Go компилируется в нативный исполняемый файл, а модель конкурентности на основе горутин
позволяет обслуживать множество одновременных соединений без выделения отдельного потока
операционной системы на каждый запрос \cite{saioc2023goroutine}.
Требование по времени ответа обеспечивается за счeт использования пула соединений
драйвера \texttt{pgx/v5}, кэширования расстояний и перекладывания вычислений
матрицы расстояний на клиентскую сторону (Яндекс Карты MultiRoute),
что снижает нагрузку на сервер при оптимизации.
Устойчивость к повторным запросам к внешним API обеспечивается таблицей
\texttt{api\_distance\_cache}: при наличии ранее вычисленного расстояния
между парой точек обращение к OSRM не производится.
Контейнеризация на основе Docker Compose обеспечивает воспроизводимость
развeртывания и независимость от конфигурации целевого сервера.
Поддержка современных браузеров обеспечивается выбором клиентского стека
(React~18 + Vite~6 + TypeScript): согласно официальной документации,
React~18 поддерживает все браузеры, совместимые с ES2015~\cite{react},
а Vite~6 нацеливается на браузеры с поддержкой нативных ES-модулей~\cite{vite}
(Chrome~89+, Firefox~89+, Safari~15+, Edge~89+).

Сравнение запланированных требований с результатами реализации представлено
в таблице~\ref{tab:requirements_check}.

\begin{table}[H]
  \caption{Соответствие реализованной системы поставленным требованиям}
  \label{tab:requirements_check}
  \begin{tabular}{|p{0.55\textwidth}|p{0.35\textwidth}|}
    \hline
    \textbf{Требование}                                               & \textbf{Результат}                                                                                \\ \hline
    Создание, редактирование и удаление задач                         & Реализовано                                                                                       \\ \hline
    Хранение географических координат                                 & Реализовано                                                                                       \\ \hline
    Получение матрицы расстояний через внешние API                    & Реализовано (OSRM + кэш)                                                                          \\ \hline
    Вычисление оптимальной последовательности с учeтом временных окон & Реализовано                                                                                       \\ \hline
    Расчeт времени прибытия, ожидания и завершения                    & Реализовано                                                                                       \\ \hline
    Визуализация маршрута на карте                                    & Реализовано (Яндекс Карты)                                                                        \\ \hline
    Экспорт задач в форматах CSV, XLSX и JSON                          & Реализовано                                                                                       \\ \hline
    Импорт задач из CSV и XLSX                                        & Реализовано                                                                                       \\ \hline
    Дополнительные ограничения маршрута (начальная/конечная точка, порядок) & Реализовано                                                                                   \\ \hline
    Выбор режима транспорта                                           & Реализовано (авто, пешком, общественный транспорт)                                                 \\ \hline
    Обнаружение конфликтов временных окон                             & Реализовано                                                                                       \\ \hline
    Время ответа сервера не более 2~секунд                            & Выполнено                                                                                         \\ \hline
    Поддержка современных браузеров                                   & Выполнено                                                                                         \\ \hline
    Устойчивость к повторным запросам API за счeт кэширования         & Выполнено                                                                                         \\ \hline
    Масштабируемость до 100 одновременных пользователей               & Подтверждено нагрузочным тестированием: при 100~VU p95 = 1370~мс $<$ 2000~мс, доля ошибок~— 0\,\% \\ \hline
  \end{tabular}
\end{table}

Таким образом, все функциональные и нефункциональные требования выполнены в полном объeме,
что подтверждает успешность реализации разработанного приложения.


\section{Аспекты отказоустойчивости, надeжности и производительности}

\subsection{Устойчивость к отказам внешних зависимостей}

Одним из ключевых рисков систем, интегрирующих внешние сервисы, является
деградация или временная недоступность стороннего API.
В разработанном приложении этот риск снижается на двух уровнях.

На уровне базы данных реализована таблица \texttt{api\_distance\_cache},
хранящая ранее полученные значения расстояния и времени перемещения
между парами географических точек.
При наличии записи в кэше сервер не выполняет повторного обращения к OSRM.
Это снижает зависимость от доступности внешнего маршрутизатора и уменьшает
суммарную задержку при повторных оптимизациях маршрутов с теми же точками.

На уровне пользовательского интерфейса загрузка Яндекс Карт реализована
через утилиту \texttt{loadYandexMaps()}, которая инжектирует \texttt{<script>}-тег
динамически при первом обращении и возвращает один и тот же Promise при
последующих вызовах (одиночный экземпляр).
Такой подход предотвращает дублирующую загрузку скрипта при одновременном
обращении нескольких компонентов и исключает состояние гонки при инициализации API.

\subsection{Надeжность аутентификации}

Аутентификация реализована по схеме JWT~\cite{rfc7519} с двумя токенами:
краткосрочным \texttt{access}-токеном (время жизни~— 15~минут) и
долгосрочным \texttt{refresh}-токеном (время жизни~— 30~дней).
Токены обновления хранятся в базе данных в виде хэша \texttt{bcrypt}~\cite{provos1999bcrypt},
что исключает их компрометацию при несанкционированном доступе к СУБД.

На клиентской стороне реализован механизм автоматического обновления токена
при получении ответа с кодом 401.
Для исключения множественных одновременных запросов на обновление применяется
дедупликация: при первом истечении токена запускается единственный запрос
обновления, остальные запросы ожидают его завершения через общий Promise.
После успешного обновления исходный запрос повторяется автоматически —
пользователь не замечает прерывания сессии.

Сессия пользователя сохраняется в \texttt{localStorage} под ключом
\texttt{planner.auth.session}.
При закрытии и повторном открытии браузера приложение автоматически
восстанавливает состояние аутентификации без повторного ввода учeтных данных.

Диаграммы последовательности, отражающие три сценария
взаимодействия компонентов системы аутентификации, представлены
на рисунках~\ref{fig:auth_sequence_login}, \ref{fig:auth_sequence_valid}
и~\ref{fig:auth_sequence_refresh}.

\begin{figure}[H]
  \centering
  \includegraphics[width=0.95\textwidth]{auth_sequence_login.png}
  \caption{Диаграмма последовательности: сценарий входа в систему}
  \label{fig:auth_sequence_login}
\end{figure}

\begin{figure}[H]
  \centering
  \includegraphics[width=0.95\textwidth]{auth_sequence_valid.png}
  \caption{Диаграмма последовательности: запрос с валидным токеном}
  \label{fig:auth_sequence_valid}
\end{figure}

\begin{figure}[H]
  \centering
  \includegraphics[width=0.95\textwidth]{auth_sequence_refresh.png}
  \caption{Диаграмма последовательности: автоматическое обновление токена при ответе 401}
  \label{fig:auth_sequence_refresh}
\end{figure}

\subsection{Надeжность инфраструктуры}

При развeртывании с использованием Docker Compose~\cite{docker} зависимость между контейнерами
управляется через механизм \texttt{healthcheck}: сервис применения миграций и
API-сервер запускаются только после того, как контейнер базы данных PostgreSQL
подтвердит готовность к приeму соединений командой \texttt{pg\_isready}.
Это исключает ситуацию, при которой API-сервер стартует до завершения инициализации СУБД.

Схема базы данных управляется инструментом \texttt{migrate/migrate},
запускаемым автоматически при каждом старте стека.
Применение миграций является идемпотентным: повторный запуск \texttt{migrate up}
при уже актуальной схеме не вносит изменений.
Это позволяет безопасно перезапускать стек без ручного контроля состояния миграций.

\subsection{Производительность}

Производительность серверной части обеспечивается несколькими механизмами.
Язык Go компилируется в нативный исполняемый файл, что устраняет накладные расходы
интерпретируемых сред.
Драйвер \texttt{pgx/v5} использует пул соединений к PostgreSQL,
исключая создание нового соединения на каждый HTTP-запрос.
HTTP-фреймворк Gin применяет \texttt{httprouter} с маршрутизацией на дереве-префикса
(radix tree), обеспечивая постоянное время поиска маршрута вне зависимости
от числа зарегистрированных эндпоинтов.

Вычисления матрицы расстояний при оптимизации маршрута выполняются
на стороне клиента посредством Яндекс Карт API,
что разгружает сервер от вычислительно eмких операций и
позволяет обслуживать большее число одновременных пользователей
без изменения аппаратной конфигурации.

Требование нефункциональных характеристик об устойчивости к ограничениям
внешних API выполнено за счeт кэша расстояний:
при повторных оптимизациях маршрутов с теми же парами точек обращения
к OSRM и Яндекс Картам не производятся.

\subsection{Производительность алгоритма оптимизации}

Для оценки производительности алгоритма NNH-TW использован встроенный
инструмент бенчмаркинга языка Go (\texttt{testing.Benchmark}).
Измерения проводились на рабочей станции с процессором Intel Core~i5-12400
(тактовая частота~2,5~ГГц, 6~ядер) и~16~ГБ оперативной памяти
под управлением операционной системы macOS~15.
Для каждого значения~$n$ выполнялось не менее~1000~итераций; в таблице
приведено среднее время одного вызова.
Результаты представлены в таблице~\ref{tab:algo_perf}.

\begin{table}[H]
  \caption{Производительность алгоритма NNH-TW в зависимости от числа задач}
  \label{tab:algo_perf}
  \begin{tabularx}{\linewidth}{|>{\centering\arraybackslash}p{2.2cm}|>{\centering\arraybackslash}p{4.8cm}|>{\centering\arraybackslash}p{3.2cm}|>{\centering\arraybackslash}X|}
    \hline
    \textbf{Число задач ($n$)} & \textbf{Запросов к API матрицы ($n(n-1)$)} & \textbf{Время NNH-TW (Go)} & \textbf{Узкое место}         \\ \hline
    5                          & 20                                         & $< 0{,}01$~мс              & HTTP-запросы к Яндекс Картам \\ \hline
    10                         & 90                                         & $< 0{,}1$~мс               & HTTP-запросы к Яндекс Картам \\ \hline
    20                         & 380                                        & $< 1$~мс                   & HTTP-запросы к Яндекс Картам \\ \hline
    30                         & 870                                        & $< 1$~мс                   & HTTP-запросы к Яндекс Картам \\ \hline
  \end{tabularx}
\end{table}

Из таблицы~\ref{tab:algo_perf} следует, что узким местом системы является
не сам алгоритм оптимизации (сложность $O(n^3)$, выполнение за 1--2~миллисекунды),
а получение матрицы расстояний от внешнего API.
При $n = 20$ задачах число параллельных запросов к Яндекс Картам составляет
$20 \times 19 = 380$, однако эти запросы выполняются параллельно на стороне клиента
через \texttt{Promise.all}, что сводит суммарную задержку к времени
наиболее медленного из них; доминирующая задержка определяется
сетевыми условиями и временем ответа внешнего API.

Таблица \texttt{api\_distance\_cache} снижает нагрузку на внешний API при
повторных оптимизациях: если пара точек уже была обработана ранее,
расстояние возвращается из кэша без обращения к OSRM.

\subsection{Нагрузочное тестирование}

Для проверки нефункционального требования о масштабируемости до 100~одновременных
пользователей проведено нагрузочное тестирование REST API с использованием
инструмента k6~\cite{k6}.
Тест воспроизводил типовой сценарий работы пользователя: аутентификация,
получение списка задач (\texttt{GET /api/tasks}) и создание новой задачи
(\texttt{POST /api/tasks}).
Каждый сценарий выполнялся в течение 60~секунд при фиксированном числе
виртуальных пользователей (VU).
Тестовый стенд: ноутбук с процессором Apple M2 (8~ядер), 16~ГБ оперативной памяти,
серверный стек запущен локально через Docker~Compose.
Критерии успешности: 95-й процентиль времени ответа менее 2000~мс,
доля ошибочных HTTP-ответов менее 1\,\%.

Результаты нагрузочного тестирования приведены в таблице~\ref{tab:load_test}.

\begin{table}[H]
  \caption{Результаты нагрузочного тестирования REST API}
  \label{tab:load_test}
  \begin{tabular}{|c|c|c|c|c|c|}
    \hline
    \textbf{VU} & \textbf{Запросов} & \textbf{avg, мс} & \textbf{p90, мс} & \textbf{p95, мс} & \textbf{Ошибки} \\ \hline
    10          & 1500              & 69               & 133              & 159              & 0\,\%           \\ \hline
    50          & 5484              & 219              & 475              & 577              & 0\,\%           \\ \hline
    100         & 5799              & 720              & 1250             & 1370             & 0\,\%           \\ \hline
  \end{tabular}
\end{table}

Все три сценария выдержали установленные критерии успешности: ни один запрос
не завершился ошибкой, а 95-й процентиль времени ответа при 100~одновременных
пользователях составил 1370~мс, что не превышает допустимого порога в 2000~мс.
Таким образом, нефункциональное требование масштабируемости подтверждено экспериментально.

\subsection{Вывод}

Применeнные технические решения обеспечивают устойчивость приложения
к отказам внешних зависимостей, надeжную аутентификацию с автоматическим
обновлением сессии и производительность, достаточную для обслуживания
целевой аудитории системы.
Нагрузочное тестирование подтвердило выполнение требования масштабируемости:
REST API выдерживает 100 одновременных пользователей без ошибок и с временем
ответа (p95) в пределах установленного норматива.


\section{Внедрение и сопровождение программного продукта}

\subsection{Модель распространения}

Разработанное приложение является веб-приложением с клиент-серверной
архитектурой и не требует установки специализированного программного обеспечения
на рабочее место конечного пользователя.
Пользователь взаимодействует с системой через стандартный веб-браузер.

Серверная часть распространяется в виде исходного кода и \texttt{Dockerfile},
что позволяет развернуть приложение на любой платформе с поддержкой Docker
(Linux, macOS, Windows).
Клиентская часть собирается в набор статических файлов (HTML, CSS, JavaScript)
командой \texttt{pnpm build} и публикуется на любом веб-сервере
(Nginx, Apache, Caddy и др.) без дополнительных зависимостей времени исполнения.

\subsection{Процедура первоначального развeртывания}

Рекомендуемым способом развeртывания является Docker Compose.
Подробные инструкции по установке, конфигурированию и типичным затруднениям
приведены в разделе~<<Руководство пользователя и администратора>> второй главы.
Перед переходом в продуктивную эксплуатацию необходимо обязательно изменить
учeтные данные СУБД, установить криптографически стойкое значение переменной
\texttt{APP\_JWT\_SECRET} и задать корректный адрес клиентской части в
переменной \texttt{APP\_CORS\_ORIGINS}.

\subsection{Сопровождение и обновление}

Обновление, резервное копирование и мониторинг выполняются стандартными средствами Docker.
Процедуры обновления (команда \texttt{docker compose up --build}), создания резервных копий
(\texttt{pg\_dump}) и контроля состояния контейнеров подробно описаны в руководстве
администратора второй главы.

\subsection{Ограничения и направления развития}

В текущей реализации хранение географических координат задач в базе данных
осуществляется в соответствии с условиями бесплатного тарифа Яндекс Карт,
которые допускают это для учебных и некоммерческих проектов.
Для продуктивного использования в коммерческой среде потребуется переход
на коммерческую лицензию API.

Суточный лимит бесплатного тарифа составляет 1~000~запросов к геокодеру
и 25~000~запросов к JavaScript API.
При превышении указанных лимитов рекомендуется рассмотреть переход
на коммерческий тариф либо интеграцию с альтернативным геокодером.

Среди перспективных направлений развития системы можно выделить:

\begin{itemize}
  \item реализацию фоновых уведомлений о наступлении временного окна следующей задачи;
  \item поддержку командной работы с общим пулом задач;
  \item интеграцию с внешними календарями (Google Calendar, Яндекс Календарь);
  \item расширение алгоритма оптимизации более мощными метаэвристиками
        (муравьиный алгоритм, генетический алгоритм) для повышения качества
        решений при большом числе задач;
  \item учeт реального дорожного трафика при построении матрицы расстояний.
\end{itemize}

\subsection{Вывод}

Разработанный программный продукт готов к развeртыванию в продуктивной среде.
Контейнеризация на основе Docker обеспечивает воспроизводимость установки,
автоматическое управление схемой базы данных и независимость от конфигурации
целевой платформы.
Процедуры резервного копирования, мониторинга и обновления не требуют
специализированных инструментов и выполняются стандартными средствами Docker.

% В рамках ВКР удалось реализовать следующие ключевые функции:

% \begin{enumerate}
%   \item список
% \end{enumerate}

% Пример изображения и ссылки на него \ref{fig:CVE}.

% \begin{figure}[H]
%   \centering
%   \includegraphics[width=0.9\textwidth]{CVE-2025-29927.png}
%   \captionof{figure}{Обработка запросов на переход без заголовка}
%   \label{fig:CVE}
% \end{figure}

% Пример кода и ссылка на него \ref{listing:CSRF}.

% \begin{listing}[H]
%   \caption{Подключение CSRF-токена}
%   \label{listing:CSRF}
%   \begin{mdframed}
%     \begin{minted}{typescript}
% import cookieParser from 'cookie-parser';
% import csurf from 'csurf';
% async function bootstrap() {
%   const app = await NestFactory.create(AppModule);
%   app.use(cookieParser());
%   app.use(
%     csurf({
%       cookie: {
%         httpOnly: true, // токен CSRF недоступен через JS
%         secure: true,
%         sameSite: 'strict',
%       },
%     }),
%   );
%   app.useGlobalFilters(new HttpExceptionFilter());
%   app.useGlobalPipes(new ValidationPipe({ whitelist: true }));
%   app.enableCors({origin: process.env.CORS, credentials: true});
%   await app.listen(process.env.PORT ?? 3000);
% }
% \end{minted}
%   \end{mdframed}
% \end{listing}

% \section{Вывод}

% Вывод по главе
